\chapter{Differenz und relatives Komplement}

\section{Mengendifferenz}

\begin{definition}[Differenz]
\[
 A\setminus B=\{x\in A\mid x\notin B\}.
\]
Damit gilt $x\in A\setminus B\Longleftrightarrow
x\in A\land x\notin B$.
\end{definition}

\begin{theorem}[Grundgesetze und Kennzeichen]
\[
\begin{gathered}
 A\setminus B\subseteq A,\qquad (A\setminus B)\cap B=\varnothing,\\
 A\setminus\varnothing=A,\qquad
 \varnothing\setminus A=\varnothing,\qquad A\setminus A=\varnothing,\\
 A\setminus B=\varnothing\Longleftrightarrow A\subseteq B,\qquad
 A\setminus B=A\Longleftrightarrow A\cap B=\varnothing.
\end{gathered}
\]
\end{theorem}

\begin{theorem}[Inklusionskriterium und Mengenzerlegung]
\[
\begin{gathered}
 C\subseteq A\setminus B
 \Longleftrightarrow C\subseteq A\ \land\ C\cap B=\varnothing,\\
 A=(A\cap B)\cup(A\setminus B),\qquad
 (A\cap B)\cap(A\setminus B)=\varnothing.
\end{gathered}
\]
\end{theorem}

\begin{theorem}[Inklusionen umformen und getrennt prüfen]
Für beliebige Mengen $A,B,C$ gelten
\[
\begin{aligned}
 A\setminus B\subseteq C
   &\Longleftrightarrow A\subseteq B\cup C,\\
 A\cap B\subseteq C
   &\Longleftrightarrow A\cap(B\setminus C)=\varnothing,\\
 A\subseteq B
   &\Longleftrightarrow
      \bigl(A\cap C\subseteq B\cap C\ \land\
            A\setminus C\subseteq B\setminus C\bigr).
\end{aligned}
\]
\end{theorem}

\begin{theorem}[Verhalten bei Inklusionen]
Aus $A\subseteq C$ folgt $A\setminus B\subseteq C\setminus B$.
Aus $B\subseteq C$ folgt $A\setminus C\subseteq A\setminus B$.
Insbesondere gilt
\[
 A\subseteq C,\ D\subseteq B
       \quad\Longrightarrow\quad A\setminus B\subseteq C\setminus D.
\]
\end{theorem}

\begin{theorem}[Differenz, Vereinigung und Schnitt]
\[
\begin{aligned}
 A\setminus(B\cup C)&=(A\setminus B)\cap(A\setminus C),\\
 A\setminus(B\cap C)&=(A\setminus B)\cup(A\setminus C),\\
 (A\cup B)\setminus C&=(A\setminus C)\cup(B\setminus C),\\
 (A\cap B)\setminus C&=(A\setminus C)\cap(B\setminus C),\\
 (A\setminus B)\setminus C&=A\setminus(B\cup C),\\
 A\setminus(B\setminus C)&=(A\setminus B)\cup(A\cap C),\\
 A\setminus(A\setminus B)&=A\cap B.
\end{aligned}
\]
\end{theorem}

\begin{theorem}[Existenz und Eindeutigkeit einer disjunkten Ergänzung]
Für beliebige Mengen $A,B$ gilt
\[
 \exists C\,(B\cup C=A\land B\cap C=\varnothing)
       \Longleftrightarrow B\subseteq A.
\]
Im Existenzfall ist $C$ eindeutig und gleich $A\setminus B$.
\end{theorem}

\section{Allgemeine Mengengleichungen}

Die folgenden Kriterien fassen Folgerungen der elementaren
Mengenoperationen zusammen. Jede Äquivalenz gilt für beliebige Mengen;
die Bedingungen auf der rechten Seite beschreiben genau, wann die
linke Gleichung erfüllt ist. \cite{hammack}

\begin{theorem}[Lösungen von Vereinigungs- und Schnittgleichungen]
Für beliebige Mengen $A,B,X$ gelten
\[
\begin{aligned}
 A\cup X=B
   &\Longleftrightarrow
     A\subseteq B\ \land\ B\setminus A\subseteq X\subseteq B,\\
 A\cap X=B
   &\Longleftrightarrow
     B\subseteq A\ \land\ B\subseteq X\ \land\
     (A\setminus B)\cap X=\varnothing.
\end{aligned}
\]
Die erste Gleichung besitzt genau dann eine Lösung, wenn
$A\subseteq B$; die zweite genau dann, wenn $B\subseteq A$.
\end{theorem}

\begin{theorem}[Lösungen von Differenzgleichungen]
Für beliebige Mengen $A,B,X$ gelten
\[
\begin{aligned}
 A\setminus X=B
   &\Longleftrightarrow
     B\subseteq A\ \land\ A\setminus B\subseteq X\ \land\
     B\cap X=\varnothing,\\
 X\setminus A=B
   &\Longleftrightarrow
     A\cap B=\varnothing\ \land\ B\subseteq X\subseteq A\cup B.
\end{aligned}
\]
Die erste Gleichung besitzt genau dann eine Lösung, wenn
$B\subseteq A$; die zweite genau dann, wenn $A\cap B=\varnothing$.
\end{theorem}

\begin{theorem}[Vorgegebener Schnitt und vorgegebene Vereinigung]
Für beliebige Mengen $A,B,C,X$ gilt
\[
\begin{gathered}
 A\cap X=B\ \land\ A\cup X=C\\
 \Longleftrightarrow\\
 B\subseteq A\subseteq C\ \land\ X=B\cup(C\setminus A).
\end{gathered}
\]
Das Gleichungssystem besitzt genau dann eine Lösung, wenn
$B\subseteq A\subseteq C$; in diesem Fall ist sie eindeutig.
\end{theorem}

\section{Komplement innerhalb einer festen Menge}

\begin{definition}[Relatives Komplement]
Bei einer fest gewählten Menge $U$ und $A\subseteq U$ heißt
$U\setminus A$ das \emph{Komplement von $A$ in $U$}.
Ein Komplement ohne Bezugsmenge wird nicht verwendet.
\end{definition}

\begin{theorem}[Doppelte Differenz und Inklusion im Grundbereich]
Für beliebige Mengen $U,A,B$ gelten
\[
\begin{gathered}
 U\setminus(U\setminus A)=U\cap A,\\
 A\cup(U\setminus A)=U\cup A,\qquad
 A\cap(U\setminus A)=\varnothing,\\
 U\cap A\subseteq U\cap B
       \Longleftrightarrow U\setminus B\subseteq U\setminus A,\\
 U\cap A=U\cap B
       \Longleftrightarrow U\setminus A=U\setminus B.
\end{gathered}
\]
Für $A,B\subseteq U$ vereinfachen sich $U\cap A$ zu $A$,
$U\cap B$ zu $B$ und $U\cup A$ zu $U$.
\end{theorem}

\begin{theorem}[De Morgansche Gesetze für relative Komplemente]
Für beliebige Mengen $U,A,B$ gelten
\[
\begin{aligned}
 U\setminus(A\cup B)&=(U\setminus A)\cap(U\setminus B),\\
 U\setminus(A\cap B)&=(U\setminus A)\cup(U\setminus B).
\end{aligned}
\]
\end{theorem}

\begin{theorem}[Aussonderung und Mengenoperationen]
Für eine beliebige Menge $A$ und Formeln $\varphi(x),\psi(x)$
mit festen Mengenparametern sei
$A_\varphi=\{x\in A\mid\varphi(x)\}$.
Mit derselben Schreibweise für die zusammengesetzten Formeln gilt
\[
 A_\varphi\cap A_\psi=A_{\varphi\land\psi},\qquad
 A_\varphi\cup A_\psi=A_{\varphi\lor\psi},\qquad
 A\setminus A_\varphi=A_{\neg\varphi}.
\]
\end{theorem}

\section{Symmetrische Differenz}

\begin{definition}[Symmetrische Differenz]
\[
 A\mathbin{\triangle}B=(A\setminus B)\cup(B\setminus A).
\]
Ein Element gehört genau dann dazu, wenn es in genau einer der
beiden Mengen $A,B$ liegt.
\end{definition}

\begin{theorem}[Darstellung und Gleichheitskriterium]
\[
\begin{gathered}
 A\mathbin{\triangle}B=(A\cup B)\setminus(A\cap B),\\
 A\mathbin{\triangle}B=\varnothing\Longleftrightarrow A=B,\qquad
 A\mathbin{\triangle}\varnothing=A,\qquad
 A\mathbin{\triangle}A=\varnothing.
\end{gathered}
\]
\end{theorem}

\begin{theorem}[Gleichheit innerhalb und außerhalb einer Menge]
Für beliebige Mengen $A,B,C$ gelten
\[
\begin{aligned}
 A\cup C=B\cup C
   &\Longleftrightarrow A\mathbin{\triangle}B\subseteq C\\
   &\Longleftrightarrow A\setminus C=B\setminus C,\\[.5em]
 A\cap C=B\cap C
   &\Longleftrightarrow (A\mathbin{\triangle}B)\cap C=\varnothing\\
   &\Longleftrightarrow C\setminus A=C\setminus B.
\end{aligned}
\]
\end{theorem}

\begin{theorem}[Rechengesetze der symmetrischen Differenz]
\[
\begin{aligned}
 A\mathbin{\triangle}B&=B\mathbin{\triangle}A,\\
 (A\mathbin{\triangle}B)\mathbin{\triangle}C
       &=A\mathbin{\triangle}(B\mathbin{\triangle}C),\\
 (A\mathbin{\triangle}C)\mathbin{\triangle}C&=A,\\
 A\cap(B\mathbin{\triangle}C)
       &=(A\cap B)\mathbin{\triangle}(A\cap C).
\end{aligned}
\]
Insbesondere gilt
$A\mathbin{\triangle}C=B\mathbin{\triangle}C
\Longleftrightarrow A=B$.
\cite{cameron}
\end{theorem}

\begin{theorem}[Symmetrische Differenz nach Vereinigung und Differenz]
Für beliebige Mengen $A,B,C$ gelten
\[
\begin{aligned}
 (A\cup C)\mathbin{\triangle}(B\cup C)
     &=(A\mathbin{\triangle}B)\setminus C,\\
 (A\setminus C)\mathbin{\triangle}(B\setminus C)
     &=(A\mathbin{\triangle}B)\setminus C.
\end{aligned}
\]
\end{theorem}

\begin{theorem}[Mengengleichung und Vergleich durch symmetrische Differenz]
Für beliebige Mengen $A,B,X$ gilt
\[
 A\mathbin{\triangle}X=B
       \Longleftrightarrow X=A\mathbin{\triangle}B.
\]
Insbesondere besitzt die Gleichung immer genau eine Lösung $X$.
Für beliebige Mengen $A,B,C$ gilt außerdem
\[
 A\mathbin{\triangle}C
       \subseteq(A\mathbin{\triangle}B)\cup(B\mathbin{\triangle}C).
\]
\end{theorem}
